\chapter{Scaling}

\section{What Breaks at Scale}

Chapter 26 and Chapter 29 each deferred a question to this chapter without pretending the deferral was free. Chapter 26 designed sparse, entropy-aware, locality-preserving storage, and then admitted that a world database of any real size will not fit on a single machine. Chapter 29 designed conflict resolution and social monitoring for many simultaneous agents, and then admitted that genuine reconciliation across real network latency needed its own treatment. It is worth naming, before building anything, that these are not one problem but three: storage too large for one machine's disk, computation too large for one machine's processors, and coordination too slow across one machine's network boundary to the next. This chapter addresses all three, and it does so by extending machinery this book has already built rather than introducing a separate distributed-systems architecture bolted onto the side of everything Parts II through IV established.

\section{Sharding by Patch}

The natural unit for splitting the world database across machines is, once again, Chapter 21's patches, and it is worth noting how much work this single design decision, made originally for purely mathematical reasons, continues to save this book from having to redo. Each shard, a machine or cluster of machines responsible for a portion of the world, owns a contiguous range of patches along the Hilbert-ordered layout Chapter 26.3 already established, and because Hilbert ordering preserves spatial locality, a contiguous range of patches in storage order is also a spatially coherent region of the domain, not a scattered handful of unrelated fragments. Sharding along these boundaries keeps nearby regions of the world on the same machine far more often than an arbitrary partition would, which matters directly for performance, since most proposals, per Chapter 29.2, touch a single local region and can be resolved entirely within one shard. Chapter 25's single-authoritative-copy discipline is preserved under this scheme, not abandoned: authority over any given patch still belongs to exactly one place, the discipline that made Chapter 15's account of commit meaningful in the first place. What changes is only that this single authority is now distributed across many shards rather than centralized in one machine, each shard the sole authority for its own patches.

\section{Cross-Shard Proposals}

The harder case, and the only case this chapter needs genuinely new machinery for, is a proposal whose target spans a boundary between two shards, an action affecting a region that straddles the line between one machine's patches and another's. This is rare, by construction, since Hilbert-ordered sharding keeps most locally relevant regions together, but it cannot be ruled out entirely. When it occurs, the two shards involved must coordinate before either commits: a proposal spanning the boundary is projected jointly, in the same sense Chapter 29.3 already described for two conflicting proposals within a single shard, except now the joint computation itself has to be carried out across a network connection rather than within one machine's memory. This chapter will not derive a full distributed-commit protocol; the standard two-phase pattern, tentatively reserving the affected patches on both shards, confirming joint admissibility, then committing on both sides together or rolling back on both sides together if admissibility fails, is adequate to the requirement and is not this book's own contribution to invent. What matters architecturally is that this coordination is needed only for the rare boundary-spanning case, not for the overwhelming majority of proposals that Hilbert-ordered sharding already keeps local to a single shard.

\section{Hot and Cold Patches}

Chapter 26.6 already established that patches are the natural unit of streaming, loaded into memory near an active observer and evicted when no longer needed. At the scale this chapter addresses, that streaming logic extends into a genuine multi-tier hierarchy rather than a simple memory-or-disk distinction. Hot patches, near currently active observers or recently subject to proposals, are held in fast, expensive storage, potentially replicated close to wherever the relevant agents are actually connecting from. Cold patches, far from any current observer and long since settled into stable, rarely-revisited coherence, migrate to slower, cheaper storage, sometimes on a different shard's tier entirely, without this migration meaning anything about the patch's content changes; it remains exactly as authoritative and exactly as available on request, only slower to retrieve on the rare occasion something does request it. This is not a new mechanism invented for scale. It is Chapter 26's streaming discipline, applied across a storage hierarchy with more tiers and more physical distance between them than a single machine's memory and disk ever required.

\section{Unbounded Domains}

This chapter's payoff is worth stating explicitly, because it closes a loop opened all the way back in Chapter 2. That chapter's central complaint against context windows and hidden states was that they are finite buffers, unable to distinguish a fact that fell out of range from a fact that never existed, precisely because their capacity is bounded in advance. The architecture this book has built since Chapter 6 was always meant to answer that complaint, and this chapter is where the answer becomes concrete rather than aspirational. Because sharding is elastic, new shards can be added as the domain grows, each taking authority over a fresh range of patches that did not need to be provisioned in advance, the domain X itself need never be bounded ahead of time the way a context window's token limit or a hidden state's fixed dimensionality must be. A world built this way can grow without a ceiling, not because any single buffer was made larger, but because authority over an ever-expanding domain is distributed across an ever-expanding set of independently bounded shards, none of which needs to anticipate, at the moment it comes into existence, how large the world it is part of will eventually become. This is what Chapter 1 asked for, stated only as a complaint at the time. This chapter is where the complaint receives its architectural answer.

\section{What This Chapter Has Not Solved}

It is worth being honest about the gap this chapter leaves open. Nothing here has addressed genuine hardware failure: what happens when a shard goes offline entirely, whether through replication, failover to a standby holding a synchronized copy of the same patches, or some other fault-tolerance mechanism this book has not developed. The standard answer, replicating each shard's authoritative patches across more than one physical machine and electing a new authority if the current one becomes unreachable, is well understood in distributed systems generally and is not something this book needs to reinvent, but it is also not something this chapter has actually specified. That specification belongs with the fuller Rust software architecture this book's appendices reserve for implementation-level detail, not the main text, where the goal has been the shape of the solution rather than its complete engineering.

\section{Closing Part V}

This completes the engineering pass Part V set out to perform. Chapter 25 gave the core write cycle a concrete software architecture, scheduler, world database, proposal and projection services, and renderers as clients. Chapter 26 gave that world database a real storage design. Chapter 27 gave the read side, rendering, its full treatment, including the discovery that rendering an unresolved region is secretly a write. Chapter 28 gave the write side the same treatment from an agent's perspective, culminating in a single invariant that governs every path into the world database without exception. Chapter 29 extended everything to many simultaneous agents, both the ordinary case of conflicting proposals and the harder case of patterns no single proposal reveals. This chapter has extended all of it to genuine scale, an unbounded domain distributed across an elastic set of shards. Nothing in Parts II through IV remains merely mathematical. Every operator this book introduced now has a place to live as running software.

What this engine is actually for, beyond the ability to sustain a wall that remembers its own crack, is the subject Part VI takes up next.
